\documentclass{article}
\usepackage{graphicx} % Required for inserting images

\title{Freelancing}
\author{Riccardo Parrino}
\date{May 2025}

\begin{document}

\maketitle

\section{Aspetti tecnici}
\subsection{Organizzazione del lavoro}
\subsection{Stack Tecnologico}

\section{Aspetti legali}

\subsection{Fatturazione}

\subsubsection{Testo di riferimento}
Il testo di referimento per qualsiasi dubbio o problema \'e il "Testo Unico IVA".

\subsubsection{Come emettere fattura}
I dati necessari per una corretta emissione della fattura sono: partita iva destinatario, pec destinatario, codice fiscale, email, 
dettagli del prodotto o servizi, parametri fiscali (profilo fiscale) condizioni e informaizoni di pagamento, ed eventuali coordinate bancarie alla quale inviare il saldo.


\subsubsection{Quando emettere fattura}
Le principali situazioni in cui emettere fattura sono la vendita di beni e/o la prestazione di servizi. 

Per quanto riguarda la vendita di bene, la fattura deve essere emessa entro 12 giorni dalla data di effettuazione dell'operazione (quindi alla consegna o alla spedizione del bene), trattandosi quindi di fattura immediata,
oppure entro il 15 del mese successivo, se si fa riferimento a documenti di trasporto (DDT) o altri documenti che provano la consegna.

Nel caso di prestazioni di servizio, se la fattura \'e immediata, questa va emessa entro 12 giorni dalla data in cui \'e stata effettuata la prestazione.
Se il pagamento avviene in anticipo, la fattura va emessa alla data del pagamento anche solo parziale. 

Per i soggetti obbligati (la maggior parte delle partite IVA), la fattura deve essere inviata al Sistema di Interscambio (SdI) entro i termini sopraindicati. La data di emissione coincide con la data di trasmissione allo SdI.

Se si dovesse emettere una fattura in ritardo, si rischiano sanzioni amministrative che possono andare dal 90\% al 180\% dell'IVA non documentata correttamente. 

\subsubsection{Sistema di Interscambio}
Per Sistema di Interscambio si intende il canale telematico dell'Agenzia delle Entrate per la gestione della fatturazione eletteonica.

Il Sistema di Interscambio riceve le fatture elettroniche, effettua i controlli su questa e le invia ai destinatari. 

Il Sistema di Interscambio verifica che siano presenti le informazioni minime obbligatorie previste per legge 
(estremi identificativi del fornitore e del cliente, il numero, la data della fattura, la descrizione della natura, la quantit\'a, la qualit\'a del bene ceduto o del servizio perstato, l'imponibile, l'aliquota e l'IVA),
che i valori della partita IVA del fornitore siano esistenti, e quindi presenti in Anagrafe Tributaria, che sia compilato il campo "Codice Destinatario" e che ci sia coerenza tra i valori dell'imponibile, dell'aliquota e dell'IVA.

Se fosse presente un qualche tipo di errore, il SdI invia una notifica di scarto e la fattura viene considerata non emessa. 

Se la fattura \'e considerata corretta allora viene trasmessa all'indirizzo riportato nel campo "Codice Destinatario" e "PEC Destinatario" del file.

\subsection{Come scrivere un contratto}
\subsection{Partita IVA}
\subsubsection{Apertura partita iva}
Per aprire partita iva \'e necessario inviare il modello AA9/12 o telematicamente o fisicamente all' Agenzia delle entrate. 

All'interno di questo modello \'e necessario inserire il proprio nome e cognome se si tratta di lavoratore autonomo, della denominazione dell'azienda se si tratta di un'impresa. Successivamente \'e necessario inserire il codice ATECO della propria attivit\'a, il tipo di regime fiscale (ordinario, ordinario semplificato o forfettario, supponendo che si soddisfino tutti i requisiti), il luogo in cui si svolge l'attivit\'a ed eventuali collaboratori come anche il prorpio commercialista.

Per quanto riguarda il luogo di svolgimento della propria attivit\'a \'e necessario indicare un luogo in cui si \'e facilmente reperibili ed \'e apposta eventuale targhetta di riconoscimento. Se si \'e in affitto, \'e necessario informare il proprietario di casa che si intende usare l'abitazione come ufficio.

Oltre a questi adempimenti in talune situazioni \'e necessario iscriversi all'INPS per gestione separata o artigiani/commercianti, all'INAIL in caso di rischio per la salute, alla Camera di Commercio (per artigiani e commercianti) e SCIA al Comune in caso di  attivit\'a soggetta a comunicazione di inizio.

Successivamente a questi passaggi \'e necessario avere una propria PEC (Posta Elettronica Certificata) e se non si rientra nel regime forfettario \'e obbligatoria la fatturazione elettronica tramite SDI. 

In genere i tempi di apertura sono immediati ed i costi sono ridotti ai soli servizi professionali (come il commercialista), i contributi INPS ed eventuali diritti camerali.

\subsubsection{Chiusura partita iva}
Inviare il modello AA9/12 di chiusura, in maniera analoga a quanto fatto per l'apertura. Non ci devono pi\'u essere fatture da emettere. Per almeno un anno, bisogno mantenere i rapporti con il proprio consulente, per la dichiarazione dei redditi e eventuali F24 da pagare per l'anno precedente.
\subsubsection{Codice ateco: Concluso}
Il codice ATECO (acronimo di ATtivit\'a ECOnomiche) \'e una classificazione alfanumerica usata in Italia per identificare il tipo di attivit\'a economica svolta da un'impresa o da un lavoratore autonomo. Il codice ATECO \'e fondamentale per l'apertura della partita IVA, la determinazione del regime fiscale, contributivo e per fini statistici. 

Il codice ATECO ha una struttura gerarchica XX.YY.ZZ.W, in cui XX individua la sezione, YY individua la divisione, ZZ il gruppo e W la Classe/Voce, anche se quest'ultimo non \'e sempre presente e serve a dare un ulteriore livello di dettaglio. Un esempio pratico: Codice ATECO: 74.10.21 significa 74 Altre attivit\'a professionali, scientifiche e tecniche, 10 si intende Attivit\'a di design specializzate e 21 Design grafico, in particolare per il web. Il codice ATECO si usa per l'apertura della partita IVA, per il calcolo dei contributi INPS, per la scelta del regime forfettario,, per statistiche ISTAT ed in alcuni casi per autorizzazione e licenze.

Per attivit\'a legate allo sviluppo software o legate all'informatica, la divisione di riferimento sono la 62 e la 63. 
Per una spiegazione dettagliata e precisa, fare ricerche al seguente link: https://www.istat.it/classificazione/classificazione-delle-attivita-economiche-ateco

\subsubsection{Il commercialista}
Cosa fa: si occupa della tenuta della contabilit\'a, 
predisposizione del bilancio di esercizio e della dichiarazione dei redditi, 
redazione dei libri contabili e libro unico del lavoro, 
svolgimento di perizie tecniche, ispezioni e revisioni amministrative.

La professione del commercialista pu\'o essere definita come quella del consulente che opera 
nel campo del diritto commerciale, del diritto tributario, della ragioneria e della contabilit\'a in generale.

Il diritto commerciale \'e una branca del diritto privato che regola i rapporti attinenti alla produzione e allo scambio della ricchezza. Pi\'u in particolare, regola ed ha per oggetto i contratti conclusi tra operatori economici e tra essi ed i loro i clienti privati ed anche gli atti e le attivit\'a delle societ\'a. 

Il diritto tributario \'e un settore del diritto finanziario che regolamente i tributi di ogni tipo.

La ragioneria \'e una disciplina che si occupa della rilevazinoe dei fenomeni aziendali e della loro traduzione in scritture cnotabili, allo scopo di fornire alle direzioni delle aziende 
o ad altri soggetti interessati le informazioni necessarie per esercitare le funzioni di controllo e di decisione.

Per contabilit\'a si intende al storia economica di un'azienda nel senso che conserva una traccia di tutte le operazioni commerciali.

Oltre al commercialista, per consulenza fiscale \'e possbili avvalersi del CAF, dei consulenti del lavoro, o direttamente all'Agenzia delle Entrate.

\subsection{Regime fiscale}

\subsubsection{Regime fiscale forfetario}
Come riportato sul sito dell'agenzia delle entrate, "il regime forfetario prevede rilevanti semplificazioni ai fini IVA e ai fini contabili, e consente, altresì, la determinazione forfetaria del reddito da assoggettare a un’unica imposta in sostituzione di quelle ordinariamente previste, nonché di accedere ad un regime contributivo opzionale per le imprese."
Inoltre, \'e possibile accedere al regime forfetario senza vincoli di tempo ne di et\'a, ma soltanto soddisfando i requisiti previsti.
Il regime forfetario \'e stato introdotto dalla legge di stabilit\'a 2015 e rappresenta il regime naturale delle persone fisiche che esercitano un'attivit\'a di impresa, arte o professione in forma individuale.

Per quanto riguarda reddito e tassazione, per chi aderisce al regime forfetario \'e prevista un'unica imposta sostitutiva delle imposte sui redditi, delle addizionali regionali e comunali e dell'IRAP, nella misura del 15\%. 
L'aliquota viene applicata sull'imponibile determinato dall'ammontare dei ricavi o dei compensi percepiti moltiplicato per il coefficiente di redditivit\'a, diversificato a seconda del codice ATECO che contraddistingue l'attivit\'a esercitata.

Per i codice ATECO 62-63 \'e previsto un coefficiente di redditivit\'a del 67\%.

Per chi aderisce al regime forfetario, per i primi cinque anni di attivit\'a si potrebbe avvalere anche di un ulteriore detassazione, un'imposta sostitutiva del 5\%, visti alcuni requisiti da soddisfare.

Il regime forfetario \'e di pi\'u semplice gestione rispetto al regime ordinario od ordinario semplificato, sia per fini IVA, sia per fini sulle imposte dirette. 

Per quanto riguarda i fini Iva, non deve essere addebitata l'Iva in fattura e non deve essere detratta l'iva sugli acquisti. Non si liquida l'imposta, non si versa e non si \'e obbligati a presentare la dichiarazione e la comunicazione annuale Iva. 
Non \'e necessario comunicare all'Ageniza delle entrate le operazioni rilevanti ai fini Iva (ovvero lo spesometro) né quelle effettuate nei confronti di operatori economici aventi sede, residenza o domicilio in Paesi cosiddetti black list. Chi applica il regime forfetario, inoltre, non ha l'obbligo di registrare i corrispettivi, le fatture emesse e le ricevute.

Le semplificazioni ai fini delle imposte sul reddito riguardano l'esonero dagli obblighi di registrazione e tenuta delle scritture contabili, ma mantenere e cnoservare i registri previsti da disposizioni diverse da quelle tributarie.
Non operano le ritenute alla fonte, non si subiscono le ritenute, non si applicano le ritenute alla fonte. 

Ai fini Iva, si ha l'obbligo di numerare e conservare le fatture di acquisto e le bollette doganali, certificare i corrispettivi, 
integrare le fatture per le operazioni di cui risultatno debitori di imposta con l'indicazinoe dell'aliquota e della relativa imposta, da versare entro il giorno 16 del mese successivo a quello di effettuazione delle operazioni, senza diritto alla detrazione dell'imposta relativa.

La legge di stabilit\'a, insieme alla legge di bilancio, costituisce la manovra di finanza pubblica per il triennio di riferimento
e rappresenta lo strumento principale di attuazione degli obiettivi programmatici definiti con la Decisione di finanza pubblica. (fonte: "https://www.rgs.mef.gov.it/VERSIONE-I/attivita_istituzionali/formazione_e_gestione_del_bilancio/bilancio_di_previsione/bilancio_finanziario/BF-2016_2018/LdS2016/index.html#:~:text=La\%20legge\%20di\%20stabilit\%C3\%A0\%2C\%20insieme,la\%20Decisione\%20di\%20finanza\%20pubblica.")

\subsubsection{Adempimenti regime fiscale forfetario}
Per quanto riguarda gli adempimenti generali del regime fiscale forfetario sono previsti i seguenti: apertura della Partita IVA, fatturazione, registrazione delle operazioni, adempimenti fiscali, calcolo del reddito imponibile e i contributi INPS.

Per l'apertuea della partita iva \'e necessaria la scelta del codice ateco e la comunicazione dell'inizio attivit\'a presso l'Agenzia delle Entrate.

Per la fatturazione \'e obbligatoria l'emissione di fatture elettroniche, in cui vanno inseriti tutt gli elementi della fattura ed in cui bisogno applicare la marca da bollo da 2 euro per fatture superiori a 77,47 euro.

Per la Registrazione delle Operazioni, non \'e obbligatorio tenere registri IVA o libri contabili, ma \'e utile tenere traccia degli incassi e dei pagamenti per una gestione accurata e per future verifiche fiscali. 
In ogni caso \'e necessario conservare tutte le fatture emesse e ricevute.

Per gli adempimenti fiscali, i forfettari no applicano l'IVA sulle fatture e non la detraggono sugli acquisti. E' prevista l'imposta sostitutiva per il pagamento delle imposte e dei contributi previdenziali presso le proprie casse di appartenenza.

Per le dichiarazioni e le comunicazioni, sono previste le dichiarazione dei redditi e la certificazione unica se si hanno dipendenti o collaboratori.

\subsubsection{Scadenze regime fiscale forfetario}

\begin{comment}
    tema imposte
    tema contributi
\end{comment}

I principali adempimenti di chi aderisce al regime forfetario sono il versamento delle imposte, il versamento dei contributi previdenziali, il versamento del diritto comerale annuale, 
invio delle fatture elettroniche e la dichiarazione dei redditi.

Per quanto riguarda il versamento delle imposte, ci sono due principali scadenze: il 30 giugno ed il 30 novembre.
Il 30 giugno si versa il saldo dell'imposta sostitutiva riferita all'anno precedente e il 50\% dell'acconto dell'imposta sostitutiva dell'anno in corso. I versamenti del 30 giugno possono essere rateizzati fino a 6 rati, con cadenza mensile.

L'altra scadenza \'e invece quella del 30 novembre, procedendo al versamento del restante 50\% dell'acconto per le imposte dell'anno in corso. Quest'ultimo versamento non pu\'o essere rateizzato.

Per chi apre partita IVA, durante il primo anno non dovr\'a pagare alcuna imposta sostitutiva durante il primo anno, 
ma durante il secondo anno verser\'a l'ammontare del primo anno e l'acconto del secondo anno.

I contribuenti forfetari pagano un'imposta sostitutiva di IRPEF e addizionali regionali/comunali pari al 15\%, su un imponibile determinato sulla base del totale dei ricavi e dei compensi incassati (vige il principio di cassa) applicando la percentuale di redditivit\'a,
determinata sulla base del codice ateco, e deducendone i contributi previdenziali versati nel periodo.

Per principio di cassa si intende il sistema di registrazione dei costi e dei ricavi all'interno del bilancio secondo la quale vengono registrati soltanto i costi ed i ricavi effettivamente pagati o incassati, non soltanto emessi per fattura come avviene per il principio di competenza.
Ripetendo, con il principio di cassa vengono imputate all'interno del bilancio le sole registrazioni che hanno avuto un movimentato di denaro. Esemplificando, una fattura emessa il 31 dicembre, pagata nell'anno successivo, non sar\'a registrata nel bilancio annuale.

Questo sistema ha il vantaggio di dover pagare le imposte solo dopo aver incassato i pagamenti relativi alle fatture emesse.

Oltre al versamento dell'imposta sostitutiva \'e prevista anche l'imposta di bollo sulle fatture elettroniche, ovvero:
per ogni fattura emessa con importo superiore a 77,47 euro, per la cifra di 2 euro per ciascuna fattura. Le scadenze previste sono le seguenti:
per le fatture emesse nel primo, secondo e terzo trimestre, entro il 30 novembre, mentre per le fatture trasmesse nel quarto trimestre (quindi, ottobre-dicembre) entro il 28 febbraio. 

Oltre alle imposte, il soggetto a partita IVA forfetaria \'e tenuto al versamento dei contributi previdenziali.
L'aspetto contributivo dipende dall'attivit\'a svolta dal contribuente.
L'ammontare dovuto ed i termini in cui pagarlo \'e determinato dall'attivit\'a del contribuente.

Innanzi tutto bisogna distinguere in contribuenti che hanno una propria cassa previdenziale ed in coloro che non ce l'hanno e devono quindi affidarsi alla Gestione Separata INPS.
Ed inoltre, tra coloro che non hanno una propria cassa previdenziali, gli artigiani ed i commercianti dai restanti contribuenti. 

Per chi ha una propria cassa previdenziale, come gli avvocati, gli ingegneri e gli architetti ecc., le modalit\'a e le scadenze sono definite dalla cassa previdenziale stessa. 

Per chi invece non ha una propria cassa previdenziale deve iscriversi alla Gestione Separa INPS, pagando ogni anno un versamento in misura percentuale sul reddito forfetizzato di ogni anno.
Le aliquote variano tra il 26,07\% per chi non ha alcuna altra copertura previdenziale ed il 24,00\% per chi \'e pensionato o \'e gi\'a iscritto ad un'altra forma di previdenza obbligatoria.

Le scadenze previste per questi versamenti sono le stesse per quelle previste dall'imposta sostitutiva, 30 giugno e 30 novembre.

Come si diceva, vi \'e un'ulteriore distinsione per Artigiani e Commercianti, secondo la quale \'e prevista un'imposta di reddito detta "minimale" di 18.415 euro, da versare entro l'anno, indipendentemente dal reddito prodotto.
In altre parole, anche se in un determinato anno non si produce alcune reddito, si \'e comunque tenuti a pagare alla Gestione separata INPS, un contributo di 18.415 euro.

Se il proprio reddito, invece, supera l'importo minimale, allora andranno versati ulteriori contributi pari al 24\% (da verificare ogni anno) sul reddito eccedente.

Le scadenze sono articolate in quattro date, ovvero: 16 maggio, 20 agosto, 16 novembre, 16 febbraio (dell'anno successivo).

Per le attivit\'a che devono essere registrate alla camera di commercio (come gli artigiani ed i commercianti) \'e previsto anche il versamento del diritto camerale annuale alla Camera di Commercio. L'importo oscilla tra 44 e 53 euro. L'import deve essere versato entro il 30 giugno di ogni anno.

Dal primo gennaio 2024 anche i contribuenti in regime forfetario devono emettere fatture elettroniche, per beni (al momento della consegna o spedizione, incasso del corrispettivo se precedente la consegna o la spedizione)
e servizi (al momento dell'incasso del corrispettivo). In caso di fatturazione differita di beni o servizi, la fattura deve essere emessa e trasmessa entro il giorno 15 del mese successivo.

Infine, per chi applica il regime forfettario \'e prevista la repsentazione della dischiarazione dei redditi annuale, prevista entro il 30 settembre di ogni anno. 

\subsubsection{Come si pagano le imposte ed i contributi}
Sia le imposte che i contributi possono essere pagate tramite modello F24. 
Il modello F24 pu\'o essere pagato online tramite home banking, tramite commercialista o tramite banca o posta.

In alternativa al modello F24 esistono comunque diversi servizi telematici come FiscoOnline ed Entratel. Spesso comunque, affidandosi ad un commercialista, questo si occupa della trasmissione e del pagamento di contributi e imposte.

\subsubsection{Modello F24}
Il modello F24 \'e il modulo ufficiale usato in Italia per pagare tutte le imposte, tasse e contributi dovuti allo Stato, all'INPS, alle Regioni e ad altri enti pubblici.

Serve per pagare IRPEF, Addizionale comunale IRPEF, Imposta comunale sugli immobili (ICI), IMU, Tarsu, Tosap/cosap, Interessi pagamenti rateali, Addizionale regionale IRPEF, Tares, Tasi.

In maniera molto riassuntiva \'e composto da diverse sezioni, quali i dati anagrafici del contribuente, i codici tributo, gli importi a debito o a credito ed il periodo di riferimento.

Il modello F24 pu\'o essere pagato tramite home banking o Fisconline, tramite commercialista oppure in banca o alla posta.

Esistono diversi tipi di modello F24 tra cui ordinario utilizzato per la maggior parte dei pagamenti, il semplificato utilizzato per tributi locali come IMU, TARI, ecc. ed il precompilato fornito solitamente dal commercialista.

Per il pagamento di un determinato tributo, almeno in maniera minima, sono necessari i dati del contribuente (codice fiscale, nome, cognome, indirizzo, comune, provincia, CAP) 
ed i dati della sezione INPS (codice ente, causale contributo, codice tributo, anno di imposta, importo a debito).   

\subsubsection{Esempi di compilazione modello F24}
Per le imposte, supponendone il normale pagamento,
queste possono essere IRPEF saldo, IRPEF acconto prima rata, IRPEF acconto seconda rata o acconto in unica soluzione.

Per il pagamento IRPEF saldo bisogna compilare la sezione erario, ed inserire i seguenti dati: 
Codice Tributo(4001), 
rateazione (n.rata - numero di rate in cui si vuole frazionare l'importo totale, e.g. '0107'), 
anno di riferimento, 
importi a debito versati, 
importi a credito compensati (NON COMPILARE),
TOTALE A (Somma degli importi a debito indicati nella Sezione Erario),
TOTALE B (Somma degli import a credito indicati nella Sezione Erario, non compilare se non presenti importi a credito)
+/- (indicare il segno se positivo o negativo),
SALDO (A - B) (indicare il saldo TOTALE A - TOTALE B),
codice ufficio (NON COMPILARE),
codice atto (indicare il Codice dell'Atto oggetto di definizione)

Esempio imposta di bollo:

Esempi contributi:

\subsubsection{Regime fiscale ordinario}
Nel caso del regime fiscale ordinario sono necessari molti pi\'u adempimenti rispetto al regime ordinario. In parte dipendono se si tratta di un libero professionista o di una societ\'a, ma vediamo in generali gli adempimenti generali da dover svolgere.

Per quanto riguarda la contabilit\'a sono previste:
la registrazione delle fatture attive e passive,
le scrittture contabili obbligatorie,
la redazione di bilancio d'esercizio (per le societ\'a)

Per quanto riguarda gli adempimenti IVA sono previste:
la liquidazione periodiche IVA,
la dichiarazione IVA annuale,
la comunciazione LIPE e 
la fatturazione elettronica obbligatoria.

Per quanto riguarda le imposte sono previste:
il versamento di IRPEF o IRES (per persone fisiche o ditte individuali e societ\'a rispettivamente),
le addizionali comunali e regionali all'IRPEF,
l'IRAP, ovvero l'Imposta Regionale sulle Attivit\'a Produttive e calcolata su base imponibile diversa da IRPEF/IRES.

Le ritenute d'acconto

Per quanto riguarda le comunicazioni fiscali sono previste:
la dichiarazione dei redditi, 
la dichiarazione IRAP,
lo spesometro / esterometro (per operazioni estere) e 
le comunicazioni INTRASTAT.

Per quanto riguarda i contributi previdenziali sono previste:
l'iscrizione alla propria cassa previdenziale e
se si hanno dei dipendenti, bisogna adempiere ai vari obblighi verso INAIL ed INPS dipendenti.

Scendendo nel dettaglio di ognuno di essi:

Per quanto riguarda i libri contabili da mantenere obbligatoriamente:
libro giornale, libro degli inventari, scritture ausiliarie (conti di mastro, scritture di magazzino), 
registro dei beni ammortizzabili, registri previsti dall normativa Iva, libri sociali obbligatori previsti ai fini civilistici.

\subsection{Regime fiscale ordinario semplificato}
\subsubsection{Requisiti}
Vi possono irentrare tutte le imprese individuali e le societ\'a di persone se i loro ricavi nell'arco di un anno solare 
non superano i seguenti limiti: 500.000 euro per le prestazioni di servizi e 800.000 euro per tutte le altre attivit\'a.
Per il professionisti non vige alcun limite.

Si accede al regime semplificato indicando un volume d'affari che non supera le soglie di ricavi inferiore a quelle indicate al momento dell'apertura della partita IVA.

\subsubsection{Fiscalit\'a}
La tassazione segue le stesse regole previste il regime ordinario. Per le persone fisiche, la tassazione viene determinata in modo progressivo, bassandosi sulle aliquote IRPEF.

\subsubsection{Contabilit\'a}
A differenza del regime ordinario, nel semplifcato non si ha l'obbligo di alcune scritture contabili e non si ha l'obbligo di redigere il bilancio a fine anno.

Gli obblighi di contabilit\'a in questo regime prevedono: la tenuta del registro dei beni ammortizzabili, la tenute del registro IVA, la tenuta del registro degli incassi dei pagamenti e la tenuta del libro unico del lavoro. 

Anche nel regime semplificato (come in tutti gli altri regimi) \'e obbligatorio la fattura elettronica.

\subsubsection{Cassa previdenziale}
\subsubsection{Gestione Artigiani e Commercianti INPS}

\subsubsection{Cartelle esattoriali}
Le cartelle esattoriali sono atti ufficiali di riscossione meessa dall'agenzia delle Entrate per richiesere il pagamento di debiti verso lo Stato o altri enti pubblici. In genere contengono: l'importo da pagare (comprensivo di debito originario, sanzioni, interessi di mora, eventuali spese di notifica e riscossione), l'ente creditore (es. INPS, Comune, Agenzia delle Entrate, ...), le istruzioni per il pagamento, il termine per pagare o fare ricorso. 

Esistono diversi motivi per cui si potrebbero ricevere delle cartelle esattoriali, come ad esempio: Omissioni o ritardi nei versamenti fiscali (IRPED, IVA, imposta sostitutiva forfettaria), Contribuit INPS non versati, Sanzioni per errori fiscali, multe per mancata comunicazione (esterometro, CU omesse, etc.), debiti pregressi di cui si \'e a conoscenza. Se non si paga una determinata cartella esattoriale si rischia un fermo amministrativo su un auto, il pignoramento di un conto corrente o dello stipendio, o un'ipoteca su immobili. 

Ricevendo una cartella, \'e bene verificare che sia corretta, controllando che il debito sia erale e verificare che non sia caduto in prescrizione, si pu\'o rateizzare il pagamento o si pu\'o fare ricorso entro 60 giorni se la si ritiene errata. 

\subsubsection{Contenziosi: Agenzia delle Entrate}
\begin{itemize}
    \item Omissioni o errori nella fatturazione
    \item Compensi non dichiarati
    \item Superamento soglie del regime forfettario, ma non si passa al regime ordinario
    \item Deducibilità di spese non inerenti (viaggi, cene, prodotti ad uso personale non legate all'attivit\'a produttiva)
    \item Mancata iscrizione alla gestione previdenziale INPS (o errata)
    \item Errori nei versamenti di imposte e contributi
    \item Errori nei rapporti con clienti esteri (mancata iscrizione al VIES, reverse charge non applicato)
\end{itemize}

\subsubsection{Contenziosi: Clienti}
\begin{itemize}
    \item Mancato pagamento o ritardo nei pagamenti
    \item Ambiguità o assenza di contratto
    \item Scope creep (espansione non concordata delle richieste)
    \item Diritti di propriet\'a intellettuale
    \item Violazione della privacy o del GDPR
    \item Contenziosi con agenzie o intermediari
    \item Conflitti sulla qualit\'a del software
\end{itemize}

\subsection{Diritto commerciale}

L'imprenditore all'art.2082: imprenditore \0e chi esercita progessionalmente un'attitivit\' economica organizzata al fine della prodczione e dello scambio di beni e eservizi

attivit\'a economica: sopportare cost e realizzare ricavi.

Attivit\'a organizzata: presenza di capitale e elavoro.

Attivit\'a professionale: svolta in modo abitale, costante, duraturo, non occasionale anche se stagionale.

Classificazione di imprese:
impresa commerciale: imprenditore non piccolo, non agricola: industriale, trasporto per terra, acqua, mare o aria, attivit\'a bancaria, assicurativa, acquisto e vendita beni.

impresa agricola: \'e imprenditore agricolo chi esercita un'attivit\'a diretta alla coltivazione del fondo.

Per le imprese agricole: no obbligo di tenuta delle scritture contabili, no fallimento, sia piccole che grande imperse.

Piccolo imprenditore: coltivatore diretto del fondo, artifiano, piccoli commercianti. Il piccolo imprenditore non \'e soggetto a pubblicit\'a legale,
non \'e soggetto a fallimento in determinati parametri.

In particoalre, i piccoli imprenditori sono solo gli imprenditori individuali, non i collettivi.

Grande imprenditore: coloro che significativo investimento di capitale, con fallimento ed altre procedure concorsuali.

Sono soggetti all'obbligo delle iscrizione nel registro delle imprese.

Grandi o medie, legalmente non vi \'e differenza.

Classificazione dell'impresa per relazione al soggetto esercente: impresa pubblica e impresa privata.
Le prime sono imprese esercitate fagli enti pubblici, distinguiamo inoltre gli enti pubblici economici dagli enti pubblici non economici.

E le imprese private, cio\'e gestite da privati.


Registro delle imprese:
strumento di pubblicit\'a legale. serve a registrare i fatti e gi atti pi\0u sgnificativi. Si trova presso l'ufficio del registro. Chi h 

Obbligo della registrazione: imprenditori commerciaaali individuali, societ\'a commerciali, mutualistiche, enti pubblici enconomici, societ\' bancarie, assicurative, attivit\'a ausiliarie delle precedente, i consorzi. 

Tenuta informatica delle scrittture contabili: vidimizaione, numerazione e tenuta infoemarica.

Scritture contabili
documenti che contengono la situzione patrimoniale dell'impresa. Rappresentano il modo in cui sta andndo l'azienda. Rendere effeciente l'organizzazione e la gestione dell'imrprea. Tutela gli stessi imprenditori tra creditori e debitori. 
Imprenditore commerciale: ha l'obbligio di tenere il libro giornale, lirbo degli inventari (attivit\'a e passivit\'a dell'imprenditore). 
Altre scritture contabili: libro mastro, libro magazzino, libro cassa, corrsipndenze: lettere, telegrammi, fatture ricevute, fatture spedite, corretta tenuta delle scritture contabili.



Societ\'a in generale
2247 cc: societ\'a contratto tra due o pi\'u persone. Quando cresce una societ\'a  tra le parti.

Contratto consensulae, esercizio in comune di un0attivit\'a economica, conferminto di beni e serviz da parte dei soci. 
Scopo lucrativo: divisione utili tra i soci
Patrimonio sociale: beni conferiti in propriet\'a non appartengono pu\0u ai singoli soci
gesione comune dell'attivit\'a econmica.


Societ\'a di persone: societ\'a semplice, nome collettivo, accomnadita semplice
non sono persone persone giuridiche, scopo di lucro, autonomia patrimonaile imperfetta, prevalenza dell'elemento personale sul patrimonio. 


Societ\'a di capitali:
per azioni, in accomndita per azione
a responsabilit1'a limitata.
sono persone giuridiche, scopo di lucro, autonomia perfetta


Autonomia patrimoniale: 
imperfetta: patrimonio ditta non \'e separato da quello dei soci.
perfetta: patrimonio della socit\'a \'e separato da quello dei soci. 


SAPA e SRL
combinazione tra spa e sas, azionisit accomandatari e accomandanti

SRL
possiede quote di partecipazione e non azioni
capitale minimo di 10.000 euro
trarre mezzi finanziari dalle risorse di un ristretto gruppo di soci. 

srl a socio unico o unipersonale.
iscirzione nel registro delle imprese
autonomia patrimoniale

ditta individuale:
definizione e requisiti di costituzione: iscrizione al registro delle imprese, apertura di una partita IVA, autorizzazione specifiche se necessarie.
caratteristiche principali: responsabilit\'a illimitata, autonomia gestionale, fiscalit\'a semplice in base al regime scelto.
Vantaggi dell'impresa individuale: semplicit\'a burocratica, costi ridotti, flessibilit\'a operativa, possibilit\'a di agevolazioni fiscali, come il regime forfettario.
Svantaggi e rischi: responsabilit\'a patrimoniale illimitata, limitata capacit\'a di acesso al credito, crescita limitata, carico fiscale crescente. 
Se l'azienda individuale cresce notevolmente \'e necessaria ed indicata la trasformazione in un'azienda di capitali, per la separazione del patrimonio personale da quello dell'impresa. Questa operazione richiede procedure formali e approvazione di un notaio.

I diversi modi con la quale \'e possibile avviare un business sono i seguenti, per ordine di complessit\'a:
partita IVA, ditta individuale, srls, srl, Srl (Startup Innovativa), S.P.A.


\subsubsection{Ditta individuale}
Per aprire una ditta individuale \'e necessario completare alcuni pochi passaggi. Per prima cosa scegliere una ragione sociale della ditta,
Il secondo passaggio consiste nel presentare la Comunicazione Unica al Registro delle imprese, presentadosi come imprenditore commerciale o come piccolo imprenditore.
Attraverso questo secondo passaggio, sar\'a possibile adempiere a tutte le formalit\'a necessarie per costituire un'unica impresa: attribuzione Partita Iva, iscrizione al Registro delle Imprese ed eventualmente all'Albo delle Imprese Artigiane,
adempimenti INPS ai fini previdenziali, adempimenti INAIL ai fini assicurativi, e Segnalazione Certficata di Inizio Attivit\'a (SCIA) per lo Sportello Unico delle Attivit\'a Produttive (SUAP).
Sar\'a inoltre necessario avere un acasella di Posta elettronica certificata (PEC) e di un dispositivo per la Firma digitale del titolare d'impresa o del suo delegato.

I costi di apertura di un'azienda individuale sono relativamente bassi e si possono dividere in costi fissi e variabili. 
I costi fissi riguardano l'imposta di bollo, i diritti di segreteria della Camera di Commercio ed il diritto camerale. 
I costi variabili cambiano a seconda dell'attivit\'a esercitata e sono l'Attestazione di inizio attivit\'a e registrazione al SUAP, se l'attivit\'a esercitata \'e soggetta a SCIA.

Per quanto riguarda la tassazione, dipende dal regime fiscale alla quale si decide di aderire, ma supponendo una fase iniziale ed agevolata, \'e consiglabile un regime forfetario.
Secondo questo regime \'e prevista un'imposta sostitutiva al 15\% di tutte le imposte ordinario, applicata alla somma dei ricavi e dei compensi ottenuti dalla ditta all'interno di un anno, moltiplicati per il coefficienti di redditivi\'a, deciso dal codice ATECO della professione. 

A questi costi si aggiungono i contributi INPS, che seguono gli stessi principi e regole che valgono per i liberi professionisti.

\subsubsection{Libero Professionista o Ditta Individuale?}
Se si rientra nella tipologia di commerciante o artigiano allora si dovr\'a aprire partita iva come ditta individuale.

\subsection{Diritto tributario}

\subsection{Ragioneria}
Fasi di vita aziendale: organizzazione, gestione e rilevazione

rilevazione: rappresentazione dei dati che emergono dalle op. aziendali
fase tipica e fondamentale dell'azienda
alla base della rilevazione, il sistema informativo aziendale
persone, software e strumenti per la raccolta e l'analisi di queste informazioni

dare informazioni, essere disponibili, per interscambio tra le unit\'a operanti

aggregazione dei dati, quindi informazioni per il suppot delle operazioni e comunicazione con l'esterno

produrre e gestire una notevole quantit\'a di informaizoni, limitare gli errori, smistare le info in tempo reale, raccolta di info con modalit\'a prefissate

si ottimizzano le informationi: sistema dei debiti, il sistema dei crediti, il sistema dei cespiti, il sistema del personale, logistico e del magazzino, finanziario per entrate e euscite, ricerca e sviluppo, marketing ecc

offrire e ricevere informaizoni

progettato per prendere decisioni
reuqisito di tempestitvit\'a e dato da tutte le aree della socità

sottosistema contabile e non contabile
contabile: relative alle op. monetarie
inventari, scritture elementari, contabilit\'a sezionali, contabilit'a generale, contabilit\'a analitico-gestionale, pianificazione e controllo budgetario

sistema non contabile: statistiche, previsioni andamenti futuri, simulazione altrenative direzionali, analisti di mercato o di settore, reporting

rilevazione
rilevazione = scrittura
ha due finalit\'a: consuntive(rappresentare i dati) e di controllo

documenti originari: di prova, di autorizzazione, di controollo
fatrture, contratti, documenti di trasporti, gli assegni, le cambiali

classificazione delle scritture (rilevazioni)
grado di elaborazioni dati: elementari, sezionali, complesse
ordine logico: cronologiche, sistemtiche
obbligatoriet\'a: obbligatorie e facoltative
strumento di rilevazione: contabili, extracontabili

libro cassa e lo scadenziario dei debiti
libro giornale: scrittura complessa, le scritture pi\'u importanti enlla vita della societ\'a
gli stakeholder valutano l'andometno sulla basse delle scritture complesse

scritture cronologiche come il libro magazzino
scritture sistematiche

scritture obbligatorie

scritture contabili: usano ilconto come strumento di scrittura,
extracontabili invece sulla base di grafici e visualizzazioni

conto
definizione: insieme di scritture omogenee realtive ad un determinato oggetto
dare, avere
aprire / accednere un conto
addebitare
accreditare
tenere un conto
chiudere / spegnere un conto

classificazione dei conti:
funzionamento: bilaterali, unilaterali
grado di analisi oggetto: analitici, sintetici (co. ge.)
informaizoni relative all'oggeto sinottici (co. ge.) e descrittivi
forma: sezioni divise/contrapposte (co.ge), sezioni unite/accostate

contabilit\'a sezionali
prima nota / brogliaccio: scrittura elementare riferimento co.ge.(contabilit\'a generale)
mese, giorno, riferimento

contabilit\'a sezionali: settori azienda

contabilit\'a di cassa: entrate/uscite denaro, bolli e assegni, registratori di cassa, reversali/mandati, saldo: sempre positivo
contabilit\'a rapporti con le banche: c/c bancari, consistenza variazioni, estratto conto, operazioni bancarie, saldo: + o -
contabilit\'a fornitori: 
contabilit\'a clienti: documenti vendita e regolazioni
contabilit\'a IVA: 
contabilit\'a magazzino: movimento dei beni
contabilit\'a beni strumentali: registro cespiti ammortizzabili, libro inventari, registro acquisti: storia dei beni
contabilit\'a personale: libro matricoola, libro paga, registro infortuni

Obblighi contabili
adempimneti che le imprese sono tenute ad effettuare per manetenere coerenza tra libro contabile e 
normaitva civilistica
legislazione del lavoro
normativa fiscale
altri obblighi

Normativa civilistica:
art. 2082 c.c. -> imprenditore commerciale
art. 2195 c.c. -> impresa commerciale
art. 2214 c.c. -> imprenditore commerciale deve tenere il libro giornale ed il libro inventari
art. 2220 c.c. -> imp. commerciale deve tenere anche: libri contabili specifici azienda, per 10 anni atti che riguardano l'attivit\'a
art. 2083 c.c. -> piccolo imprenditore: coltivatore diretto, artigiano, lavoratore in proprio

scrittura contabili previsti da normativa a civilistica: 2114 a 2220
libro giornale, libro inventari, scritture contabili richieste dalla natura e dalle dimensioni dell'impresa (es. socit1'a di capitali: libro soci, libro verbali adunanze)

Legislazione del lavoro: libro matricola e libro paga: sostuiti dal libro unico, soggetti a bollatura da parte dell'INAIL
registro degli infortuni

Normativa fiscale: 
DPR 633/72 decreto iva, DPR 600/73 decreto accertamento imposte sui redditi
Decreto IVA: registro acquisti, registro fatture emesse, registro corrispettivi di vendita

regime contabile agevolato per contribuenti minimi

contabilit\'a ordinaria -> societ\'a capitali, individuali, di persone

finalit\'a tenuta e conservzione delle scritture contabili obbligatorie

art. 2219 c.c. contabilit\'a ordinata
art. 2220 c.c. conservazione per 10 anni

contabilit\'a generale
co. ge.
Principale strumento rilevazione: registrare fatti esterni di gestione
metodi di scritture: regole
metodo della partita doppia \'e applicato al sistema del capitale a patrimonio e del risultato economico
(ideato da Amaduzzi)

basato sul dare ed avere
variazioni finanziarie vs variazioni economiche

Partita doppia, conti finanziari ed economici

\subsection{Contabilit\'a}

Esistono diverse contabilit\'a per le aziende: Mercantile o commerciale, Industriale o tecnica.

Per mercantile o commerciale, si intende la contabilit\'a che collega l'azienda con il mondo esterno e riguarda l'acquisto di materie prime (uscite) e la vendita di prodotti o servizi (entrate);

Per industriale o tecnica, si intende la contabilit\'a che sta all'interno dell'azienda e riguarda i processi di trasformazione delle materie prime in prodotti finiti.

Da questa distinsione ne segue un'ulteriore, ovvero: la contabilit\'a generale (Co.Ge) e la contabilit\'a industriale (anche detta contabilit\'a analitica, Co.An).
La prima registra tutti i fatti amministrativi intercorsi tra l'azienda e l'ammbiente esterno. La seconda registra solo i fatti di gestione interna.

Nella Pubblica Amministrazione la metodologia viene detta della Contabilit\'a Finanziaria, e fino a pochi anni fa era della partita semplice. Da pochi anni \'e stato introdotto il metodo della partita doppia, con la denominazione contabilit\'a economica.

Ragioneria analizzare la contabilit\'a 
cnotabilit\'a: tutte le operazioni che un 'azienda svolge nella sua vita
metodo della partita doppia per realizzare ci\'ob
conti raccolti nel libro mastro. In due sezione, dare ed avere
ogni fatto di gestione va osservato sconod due aspetti, aspetto originario ed aspetto derivato
il tatolae degli importi "dare" deve essere uguale al totale importi "avere"

esistono diversi libri, libro giornale, libro mastro e libro inventari.
libro giornale: le operazioni fatte ogni giorno

libro inventari: obbli. all'inizio dell'impresa, insieme di tutte le passivit\'a dell'azienda

classificazioni di operazioni:
variazioni finanziarie e variazioni ECOnomiche
finanziarie: positive(aumento di denaro) e negative(diminuzione di denaro)

var. economiche: positive(ricave, aumento di capitale), negative(costi, diminuzione di capitale)

conto: + (dare, vfp, ven), - (avere, vfn, vep)

iva

Operazioni di acquisto
acquisto di beni
acquisto di servizi
acquisto d immobilizzazioni
acquisto sui mercati esteri

ogni operazoini va rilevata
correttta valorizzazione dell'immobile
momento di rilevazione
valori originati: costo del bene, debito di regolamento, iva a credito
fasi rilevazione: liquidazione, pagamento

Esempio di acquisto di beni
acquisto di beni
mat.prime c/acquisti +  iva a credito: dare 10.000 + 2.200
avere 12.200



\subsection{Tributi}

\subsubsection{INPS}
L’INPS gestisce la liquidazione e il pagamento delle pensioni e delle indennità di natura previdenziale e assistenziale. Le pensioni sono prestazioni previdenziali, determinate sulla base di rapporti assicurativi e finanziate con i contributi di lavoratori e aziende pubbliche e private. Invece, le prestazioni assistenziali o “a sostegno del reddito” tutelano i lavoratori che si trovano in particolari momenti di difficoltà della loro vita lavorativa e provvedono al pagamento di somme destinate a coloro che hanno redditi modesti e famiglie numerose. Per alcune di queste prestazioni l’INPS è coinvolto solo nella fase di erogazione, mentre per altre svolge tutto il procedimento di assegnazione (fonte: sito INPS). L’INPS nasce oltre centoventicinque anni fa, nel 1898 con la fondazione della Cassa Nazionale di previdenza per l'invalidità e per la vecchiaia degli operai, allo scopo di garantire i lavoratori dai rischi di invalidità, vecchiaia e morte. Con il tempo l’Istituto ha assunto un ruolo di crescente importanza fino a diventare il pilastro del sistema nazionale del welfare.

\subsubsection{INAIL}
L'INAIL (Istituto Nazionale per l'Assicurazione contro gli Infortuni sul Lavoro) \'e un ente pubblico italiano che si occupa di assicurazione contro gli infortuni sul lavoro.
In generale offre servizi per la tutela dei lavoratori in caso di danni alla salute causati dallo svolgimento dell'attivit\'a lavorativa, promuove iniziative per migliorare la sicurezza nei luoghi di lavoro, 
fornisce supporto medico, riabilitativo ed anche professionale ai lavoratori infortunati o malati ed eroaga compensazioni economiche in caso di infortunio o malattia professinoale che comporti inabilit\'a temporanea, permanente o, nei casi pi\'u gravi, la morte del lavoratore.

Per alcune attivit\'a considerate rischiose (artigianato, edilizia, industria, ecc.), l'assicurazione INAIL \'e obbligatoria (da parte dei datori di lavoro).

\subsubsection{Modello F24}
Il modello F24 \'e un modulo utilizzato per il pagamento di tasse, imposte, contributi e tributi e serve a versare importi dovuti all'Agenzia delle Entrate, all'INPS, ai Comuni o ad altri enti. 

Attraverso il modello F24 si possono pagare: imposte sui redditi (IRPEF, IRES), IVA, IMU, TASI, TARI, Contributi INPS e INAUL, Sanzioni e Interessi, Accise e tributi regionali e locali

Esistono diverse versioni del modello F24: F24 ordinario, F24 semplificato, F24 accise ed F24 ELIDE.

Il modello F24 pu\'o essere pagato Online tramite i servizi telematici dell'Agenzia delle Entrate o del proprio home banking, in banca o posta o tramite intermediari abilitati come i commercialisti.

\section{Aspetti economici}
\subsection{Modi di guadagnare}
\subsection{Quanto farsi pagare}

\section{Aspetti marketing}
\subsection{Trovare clienti}
\subsection{Personal Branding}

\section{Aspetti relazionali}
\subsection{Gestione clienti}

\end{document}
